\documentclass[11pt,a4paper]{article}
\usepackage[utf8]{inputenc}
\usepackage[T1]{fontenc}
\usepackage[french]{babel}
\usepackage{amsmath}
\usepackage{amssymb}
\usepackage{geometry}
\usepackage{enumitem}
\usepackage{xcolor}
\usepackage{titlesec}
\usepackage{fancyhdr}
\usepackage{booktabs}
\usepackage{array}
\usepackage{listings}
\usepackage{tcolorbox}
\usepackage{graphicx}

% Configuration de la page
\geometry{margin=2.5cm}
\pagestyle{fancy}
\fancyhf{}
\fancyhead[L]{\leftmark}
\fancyhead[R]{K-Means Clustering}
\fancyfoot[C]{\thepage}

% Configuration des titres
\titleformat{\section}
{\Large\bfseries\color{blue!70!black}}
{}
{0em}
{}[\titlerule]

\titleformat{\subsection}
{\large\bfseries\color{blue!50!black}}
{}
{0em}
{}

% Configuration pour le code Python
\lstset{
    language=Python,
    basicstyle=\ttfamily\small,
    keywordstyle=\color{blue}\bfseries,
    commentstyle=\color{gray}\itshape,
    stringstyle=\color{red},
    numbers=left,
    numberstyle=\tiny\color{gray},
    stepnumber=1,
    numbersep=5pt,
    backgroundcolor=\color{gray!10},
    frame=single,
    breaklines=true,
    showstringspaces=false,
    tabsize=4,
    morekeywords={np, pd, sklearn, plt, sns, KMeans, train_test_split, StandardScaler, silhouette_score, accuracy_score},
    inputencoding=utf8,
    extendedchars=true,
    literate={à}{{\`a}}1 {é}{{\'e}}1 {è}{{\`e}}1 {ê}{{\^e}}1 {ë}{{\"e}}1
             {ù}{{\`u}}1 {û}{{\^u}}1 {ü}{{\"u}}1 {ô}{{\^o}}1 {ç}{{\c c}}1
             {À}{{\`A}}1 {É}{{\'E}}1 {È}{{\`E}}1 {Ê}{{\^E}}1 {Ë}{{\"E}}1
             {Ù}{{\`U}}1 {Û}{{\^U}}1 {Ü}{{\"U}}1 {Ô}{{\^O}}1 {Ç}{{\c C}}1
}

% Configuration pour les zones importantes
\tcbuselibrary{breakable}
\newtcolorbox{codebox}{
    colback=blue!5,
    colframe=blue!50,
    breakable,
    title=Code Python
}

\newtcolorbox{notebox}{
    colback=yellow!5,
    colframe=yellow!50,
    breakable,
    title=Note
}

% Métadonnées
\title{Code K-Means - Minimum Nécessaire}
\author{AmineGR03}
\date{\today}

\begin{document}

\maketitle

\section{Code Python}

\begin{codebox}
\begin{lstlisting}
# Imports
import numpy as np
import matplotlib.pyplot as plt
from sklearn.cluster import KMeans
from sklearn.model_selection import train_test_split
from sklearn.preprocessing import StandardScaler
from sklearn.metrics import silhouette_score
from sklearn.datasets import make_blobs

# 1. Chargement du dataset
X, y_true = make_blobs(n_samples=300, centers=4, n_features=2, random_state=42)

# 2. Normalisation
scaler = StandardScaler()
X_scaled = scaler.fit_transform(X)

# 3. Division train/test
X_train, X_test, _, _ = train_test_split(X_scaled, y_true, test_size=0.2, random_state=42)

# 4. K-Means
kmeans = KMeans(n_clusters=4, random_state=42, n_init=10)
kmeans.fit(X_train)

# 5. Predictions
y_train_pred = kmeans.predict(X_train)
y_test_pred = kmeans.predict(X_test)

# 6. Evaluation
silhouette_train = silhouette_score(X_train, y_train_pred)
silhouette_test = silhouette_score(X_test, y_test_pred)
print(f"Score silhouette (train): {silhouette_train:.4f}")
print(f"Score silhouette (test): {silhouette_test:.4f}")

# 7. Visualisation
plt.figure(figsize=(10, 4))

plt.subplot(1, 2, 1)
plt.scatter(X_train[:, 0], X_train[:, 1], c=y_train_pred, cmap='viridis', alpha=0.6)
plt.scatter(kmeans.cluster_centers_[:, 0], kmeans.cluster_centers_[:, 1], 
            c='red', marker='x', s=200, linewidths=3, label='Centroides')
plt.title(f'Clustering Train (Silhouette: {silhouette_train:.3f})')
plt.legend()

plt.subplot(1, 2, 2)
plt.scatter(X_test[:, 0], X_test[:, 1], c=y_test_pred, cmap='viridis', alpha=0.6)
plt.scatter(kmeans.cluster_centers_[:, 0], kmeans.cluster_centers_[:, 1], 
            c='red', marker='x', s=200, linewidths=3, label='Centroides')
plt.title(f'Clustering Test (Silhouette: {silhouette_test:.3f})')
plt.legend()

plt.tight_layout()
plt.show()
\end{lstlisting}
\end{codebox}


\end{document}

